\chapter{Grundlagen und Stand der Technik}
\label{ch:StandTechnik}
Die Extraktion robuster und invarianter Merkmale bildet eine zentrale Grundlage für bildbasierte Regelverfahren im Visual Servoing. 
Klassische Methoden wie \ac{SIFT} oder \ac{ORB} liefern stabilisierte Featurepunkte, stoßen jedoch bei komplexen Szenen und hohen 
Echtzeitanforderungen an ihre Grenzen. Moderne Ansätze aus dem Deep Learning ermöglichen hingegen eine datengetriebene und 
segmentierungsbasierte Merkmalsextraktion, die sowohl die Lokalisierung als auch die semantische Bedeutung von Bildbereichen 
verbessern kann. ?? QUELLE ??

In diesem Kapitel werden Grundlagen sowie der aktuelle Stand der Technik beleuchtet. Zunächst erfolgt ein Vergleich klassischer und moderner Verfahren zur 
Feature-Extraktion. Anschließend ein Überblick segmentierungsbasierter Methoden wie U-Net, DeepLabV3+ und 
\ac{Mask R-CNN}, die sich durch ihre Leistungsfähigkeit bei der semantischen Segmentierung und potenziellen Integration in 
Feature-Extraktion auszeichnen~\cite{potje2024xfeat}.





\section{Klassische Verfahren der Feature-Extraktion}
\label{sec:KlassischeVerfahren}
Klassische Verfahren der Feature-Extraktion basieren auf fest definierten Algorithmen, die charakteristische Bildpunkte anhand geometrischer und struktureller Merkmale erkennen.
Diese Techniken ermöglichen eine robuste Detektion und Beschreibung von Features, insbesondere hinsichtlich Skalierung, Rotation und Beleuchtung.
Des Weiteren bilden diese Grundlage vieler etablierter Anwendungen im Bereich der Robotik und Bildverarbeitung.




\subsection{\acl{SIFT}}
Das \ac{SIFT}-Verfahren, eingeführt von~\textcite{lowe2004distinctive} im Jahr 2004, ist eines der bekanntesten klassischen Verfahren zur Merkmalsextraktion in der Computer Vision.
Es wurde entwickelt, um robuste lokale Merkmale in Bildern zu identifizieren, die invariant gegenüber Skalierung, Rotation und teilweise auch Beleuchtungsänderungen sowie perspektivischen Transformationen sind.
Der Algorithmus besteht aus mehreren aufeinanderfolgenden Schritten:
\begin{enumerate}
    \item \textbf{Skalenraumextrema-Erkennung:} Hierbei wird ein Skalenraum erstellt, indem das Bild mehrfach mit Gauß-Filtern unterschiedlicher Standardabweichung gefiltert wird. Durch Subtraktion benachbarter gefilterter Bilder entstehen Differenzen von Gaußbildern, in denen potenzielle Keypoints als lokale Extrema (Minima oder Maxima) detektiert werden.
    \item \textbf{Keypoint-Lokalisierung:} Die detektierten Kandidatenpunkte werden auf ihre Stabilität geprüft, indem beispielsweise Kantenpunkte verworfen werden, die zu sensitiv gegenüber Rauschen sind.
    \item \textbf{Orientierungszuweisung:} Jedem Keypoint wird eine Orientierung zugewiesen, die auf lokalen Gradienten basiert, wodurch das Verfahren rotationsinvariant wird.
    \item \textbf{Deskriptorgenerierung:} Um jeden Keypoint wird ein Deskriptor generiert: Dabei handelt es sich um einen Vektor, der die lokalen Gradienteninformationen in der Umgebung zusammenfasst und somit den Charakter des lokalen Bildausschnitts beschreibt. Der \ac{SIFT}-Deskriptor ist 128-dimensional und ermöglicht eine robuste Wiedererkennung von Merkmalen auch unter veränderten Bedingungen.
\end{enumerate}
Aufgrund der Robustheit des \ac{SIFT}-Verfahrens, wird dieses in vielen Anwendungen wie Objekterkennung, Bildregistrierung oder Robotik eingesetzt.
Problematisch ist jedoch, dass \ac{SIFT} im Vergleich zu modernen, lernbasierten Methoden relativ rechenintensiv und daher aufwändiger zu implementieren ist.~\cite{lowe2004distinctive}.




\subsection{\acl{SURF}}
\textcite{bay2006surf} stelte im Jahr 2006 das \ac{SURF}-Verfahren vor.
Dieses ist ein beschleunigter und optimierter Nachfolger des \ac{SIFT}-Verfahrens und wie dieses zur robusten Merkmalsextraktion in Bildern geeignet. 
Ziel des \ac{SURF}-Verfahrens ist, bei vergleichbarer Robustheit eine höhere Rechengeschwindigkeit zu erreichen, sodass das Verfahren für Echtzeitanwendungen genutzt werden kann.
Erreicht wird dies u.~a.~durch den Einsatz von Integralbildern und approximativen Filteroperationen, anstelle von rechenintensiven Gauss-Filtern.
Integralbilder ermöglichen eine schnelle Berechnung von Summen über rechteckige Bildbereiche, unabhängig von deren Größe.
Im Ablauf ähnelt der \ac{SURF}-Algorithmus dem des \ac{SIFT}-Verfahrens und besteht ebenfalls aus den Schritten Keypoint-Detektion, Orientierungsschätzung und Deskriptorgenerierung:
\begin{enumerate}
    \item \textbf{Keypoint-Detektion:} \ac{SURF} verwendet eine Hessian-Matrix-basierte Methode zur Detektion von Keypoints. Dabei werden schnelle approximative Filter (Box-Filter) eingesetzt, die auf Integralbildern basieren, um die Determinante der Hessian-Matrix effizient zu berechnen.
    \item \textbf{Orientierungszuweisung:} Basierend auf Helligkeitsgradienten wird jedem Keypoint eine Orientierung zugewiesen.
    \item \textbf{Deskriptorgenerierung:} Der Deskriptor wird durch die Analyse der Verteilung von Helligkeitsgradienten in der Umgebung des Keypoints erstellt. Resultat ist ein 64-dimensionaler Vektor (in der erweiterten Version 128-dimensional) mit Information über die lokale Struktur eines Bildbereichs. 
\end{enumerate}
Ein weiterer Unterschied zu \ac{SIFT} liegt in der Verwendung von Haar-Wavelet-Filtern zur Extraktion der Gradienteninformationen.
Durch die Nutzung von Integralbildern können diese ebenfalls schnell berechnet werden können.
Durch die Invarianz gegenüber Rotation, Skalierung und partieller Beleuchtungsveränderung eignet sich \ac{SURF} für Aufgaben wie Objektwiedererkennung, Bildregistrierung und Robotik.
Allerdings ist \ac{SURF} – im Gegensatz zu \ac{SIFT} – durch Patente eingeschränkt und nicht vollständig frei verfügbar.
Dies erschwert den Einsatz in Open-Source-Projekten~\cite{bay2006surf}. 




\subsection{\acl{ORB}}
Im Jahr 2011 präsentierten~\textcite{rublee2011orb} die \ac{ORB}-Methode.
Diese zeichnet sich durch ein effizientes und patentfreies Verfahren zur Merkmalsextraktion aus, das speziell für Echtzeitanwendungen und eingebettete Systeme entwickelt wurde.
\ac{ORB} kombiniert zwei bekannte Verfahren – \ac{FAST} zur Detektion von Keypoints sowie \ac{BRIEF} zur Beschreibung – und erweitert diese um wichtige Eigenschaften 
wie Rotationsinvarianz und Robustheit gegenüber Skalierung.
Ziel ist eine leistungsfähige Alternative zu \ac{SIFT} und \ac{SURF}, die schneller arbeitet und frei verwendbar ist.
\begin{enumerate}
    \item \textbf{Keypoint-Detektion:} \ac{ORB} nutzt den \ac{FAST}-Algorithmus, der durch die Analyse von Intensitätsunterschieden in einem Kreis um einen Pixel schnelle Eckpunkte detektiert. Diese Punkte sind jedoch weder automatisch skalen- noch rotationsinvariant. Um dies zu kompensieren, wird im \ac{ORB}-Verfahren der \ac{FAST}-Algorithmus um eine Harris-Korner-Messung erweitert, um stabilere Keypoints auszuwählen.
    \item \textbf{Orientierungszuweisung:} Anschließend wird jedem Keypoint eine Hauptorientierung zugewiesen. Hierzu wird die Intensitätsverteilungen im lokalen Bildumfeld betrachtet. Diese Orientierung wird genutzt, um die Detektion und Beschreibung der Keypoints zu rotieren und somit rotationsinvariant zu machen.
    \item \textbf{Deskriptorgenerierung:} Im dritten Schritt wird der \ac{BRIEF}-Deskriptor verwendet, der auf binären Vergleichen von Pixelintensitäten basiert. Hier wird \ac{BRIEF} zu rotated \ac{BRIEF} erweitert, indem die Pixelvergleiche entsprechend der berechneten Orientierung des Keypoints angepasst werden. Dies gewährleistet, dass die Deskriptoren ebenfalls rotationsinvariant sind. Ergebnis ist ein speicher- sowie rechenfreundlicher binärer Vektor.
\end{enumerate}
Das das \ac{ORB}-Verfahren skalierbar, echtzeitfähig und frei von Lizenzbeschränkungen ist, wird es häufig in mobilen und eingebetteten Systemen eingesetzt.
Im Bezug auf die Genauigkeit und Wiedererkennung von Merkmalen in kompexen Szenen kann \ac{ORB} nicht mit modernen, lernbasierten Verfahren mithalten~\cite{rublee2011orb}. 





\section{Moderne Verfahren der Feature-Extraktion}
\label{sec:ModerneVerfahren}
Moderne Verfahren der Feature-Extraktion nutzen tiefenlernbasierte Modelle, um Merkmale direkt aus Bilddaten zu lernen. 
Einige Ansätze bieten eine hohe Robustheit gegenüber komplexen Szenen, Beleuchtungsänderungen und Verdeckungen.
Sie ermöglichen eine präzisere und anpassungsfähigere Merkmalserkennung als klassische Methoden. ?? QUELLE ??




\subsection{SuperPoint}
SuperPoint ist ein lernbasiertes Verfahren zur Detektion und Beschreibung von lokalen Bildmerkmalen, das 2018 von~\textcite{detone2018superpoint} vorgestellt wurde. Es unterscheidet sich grundlegend von klassischen Ansätzen wie bspw.~\ac{SIFT}, \ac{SURF} oder \ac{ORB}.
Anstelle eines Algorithmuses wird die Nutzung ein neuronales Netzwerks vorgeschlagen, das sowohl Keypoints erkennt als auch die zugehörigen Deskriptoren berechnet. – Dies geschieht in einer gemeinsamen, vollkonvolutionalen Pipeline.
SuperPoint wurde insbesondere für Echtzeitanwendungen in der Robotik und Augmented Reality entwickelt und zeichnet sich durch hohe Wiederholbarkeit, Robustheit und Geschwindigkeit aus.
Das Verfahren basiert auf einem zweistufigen Trainingsansatz:
\begin{enumerate}
    \item \textbf{Pre-Training auf synthetischen Daten:} Zunächst wird ein kleines Netzwerk, genannt „MagicPoint“, auf einem synthetisch generierten Datensatz trainiert. Dieser Datensatz besteht aus einfachen geometrischen Formen (z.B. Dreiecken, Quadraten) mit künstlich hinzugefügtem Rauschen. Ziel ist es, dass das Netzwerk lernt, robuste Keypoints in diesen stark vereinfachten Szenen zu erkennen.
    \item \textbf{Self-Supervised Fine-Tuning auf realen Bildern:} Im zweiten Schritt wird das vortrainierte Netzwerk mit realen Bildern weitertrainiert. Hierzu werden Bildpaare verwendet, die durch bekannte geometrische Transformationen (Homographien) miteinander verbunden sind. Das Netzwerk lernt, konsistente Keypoints und Deskriptoren zu erzeugen, ohne dass manuelle Annotationen erforderlich sind.
\end{enumerate}
Erzeugt werden zwei Ausgaben: eine Keypoint-Heatmap und eine dichte Deskriptormatrix.
Die Heatmap zeigt, wo im Bild relevante Punkte sind, während die Deskriptoren die Bildumgebung dieser Punkte als Vektoren beschreiben.
Dadurch eignet sich SuperPoint besonders gut für Aufgaben wie \ac{SLAM}, Visual Odometry oder Structure from Motion.
Ein weiterer Vorteil von SuperPoint ist die Kombination aus hoher Genauigkeit und Echtzeitfähigkeit in Verbindung einer \ac{GPU}.
SuperPoint ist robust gegenüber Skalierung, Rotation sowie moderaten Beleuchtungsänderungen.
Des Weiteren ermöglicht das gemeinsame Lernen von Detektion und Beschreibung eine bessere Abstimmung zwischen diesen beiden Komponenten im Vergleich zu modularen Verfahren~\cite{detone2018superpoint}.

\subsection{R2D2}
Repeatable and Reliable Detector and Descriptor (R2D2) ist ein lernbasiertes Verfahren zur lokalen Merkmalsextraktion, 
das 2019 von Revaud et al. eingeführt wurde. Ziel dieses Ansatzes ist es, Bildpunkte zu identifizieren, die sowohl 
wiederholbar (also unter verschiedenen Ansichten des gleichen Objekts stabil) als auch diskriminativ (gut unterscheidbar
von anderen Punkten) sind. Im Gegensatz zu klassischen Methoden wie \ac{SIFT} oder \ac{ORB}, die fest programmierte Schritte für 
Keypoint-Detektion und Deskriptor-Bildung verwenden, setzt R2D2 auf ein neuronales Netz, das diese Eigenschaften direkt 
lernt. Das Kernprinzip von R2D2 besteht darin, ein vollkonvolutionales Netzwerk zu trainieren, das drei Karten ausgibt: 
eine Deskriptorkarte, eine Wiederholbarkeitskarte und eine Zuverlässigkeitskarte. Die Deskriptorkarte enthält Vektoren, 
die lokale Bildregionen beschreiben. Die Wiederholbarkeitskarte zeigt, wie stabil ein Punkt bei Transformationen (z.B. Rotation, Skalierung) ist, 
während die Zuverlässigkeitskarte angibt, wie gut sich ein Punkt zum Matching eignet. Die Kombination dieser drei Informationen erlaubt es, 
besonders robuste und eindeutig zuzuordnende Punkte auszuwählen. Im Training nutzt R2D2 Bildpaare mit bekannten geometrischen Transformationen 
(z.B. Homographien), um die gewünschten Eigenschaften der Keypoints zu erzwingen. Die Verlustfunktion belohnt Punkte, 
die unter diesen Transformationen wieder auftauchen (hohe Wiederholbarkeit), und bestraft Punkte, deren Deskriptoren zu 
ähnlich oder zu unspezifisch sind (niedrige Zuverlässigkeit). Dadurch entsteht ein Verfahren, das sowohl präzise als auch 
effizient ist – eine wichtige Eigenschaft für visuelle SLAM-Systeme oder Structure-from-Motion. R2D2 ist insbesondere für 
schwierige Bildbedingungen wie starke Perspektivwechsel oder Wiederholungsmuster (z.B. Zäune, Fensterreihen) geeignet. 
Durch die gleichzeitige Optimierung von Detektor- und Deskriptoreigenschaften hebt sich das Verfahren von Ansätzen wie 
SuperPoint ab, die diese Komponenten getrennt behandeln~\cite{revaud2019r2d2}. 


\subsection{LoFTR}
Local Feature TRansformer (LoFTR) ist ein modernes, lernbasiertes Verfahren zur Feature-Korrespondenzbestimmung, das 2021 
von Sun et al. vorgestellt wurde. Im Gegensatz zu klassischen Methoden wie \ac{SIFT} oder auch hybriden Deep-Learning-Ansätzen 
wie SuperPoint + SuperGlue, verzichtet LoFTR vollständig auf die explizite Detektion einzelner Keypoints. 
Stattdessen verfolgt es einen „detector-free“ Ansatz, bei dem dichte Korrespondenzen direkt zwischen zwei Bildern auf Pixel- 
oder Patch-Ebene berechnet werden. Dies ermöglicht eine deutlich robustere Zuordnung – insbesondere in texturarmen oder stark 
verzerrten Szenen, in denen klassische Verfahren oft versagen. Der Kern von LoFTR ist ein Transformer-basiertes Netzwerk, 
das Bildpaare gleichzeitig verarbeitet und über sogenannte coarse-to-fine Matching-Strategien hochpräzise Korrespondenzen 
berechnet. Zunächst werden die Bildmerkmale durch ein geteiltes CNN extrahiert, anschließend mit einem „coarse-level transformer“ 
global abgestimmt, um grobe Korrespondenzen zu finden. Danach verfeinert ein „fine-level transformer“ die Matches lokal, 
um subpixelgenaue Punktpaarungen zu erzeugen. Dabei nutzt LoFTR sowohl Selbst-Attention innerhalb eines Bildes als auch 
Cross-Attention zwischen den beiden Bildern – ähnlich wie es in der Sprachverarbeitung bei Übersetzungsmodellen bekannt ist. 
Ein bedeutender Vorteil von LoFTR ist die Fähigkeit, auch in schwierigen Szenarien wie extremen Blickwinkeländerungen, 
schwacher Textur oder Wiederholungsmustern robuste Matches zu liefern. Dies macht es besonders wertvoll für Anwendungen wie 
Structure-from-Motion, visuelle Lokalisierung und SLAM. Trotz des rechnerischen Aufwands ermöglicht eine optimierte 
Implementierung sogar den Einsatz in Echtzeitumgebungen mit GPU-Beschleunigung~\cite{sun2021loftr}. 


\section{Segmentierungsbasierte Merkmalsextraktion}
\label{sec:Merkmalsextraktion}

\subsection{Mask R-CNN}
Mask R-CNN ist ein leistungsstarkes neuronales Netzwerk zur Instanzsegmentierung, das 2017 von Kaiming He et al. vorgestellt wurde. 
Es erkennt nicht nur Objekte in Bildern und lokalisiert sie mit Bounding Boxes, sondern erzeugt zusätzlich für jedes Objekt eine binäre 
Maske, die dessen exakten Umriss pixelgenau abbildet. Die Architektur besteht aus mehreren Modulen: Ein Convolutional Neural Network (CNN)
extrahiert Bildmerkmale, ein Region Proposal Network (RPN) schlägt relevante Regionen vor, und mittels Region of Interest (RoI) Align werden diese Merkmale präzise 
auf jede Region übertragen. Für jede erkannte Region erfolgen drei parallele Vorhersagen – Objektklasse, genaue Bounding Box und Segmentierungsmaske –, 
wodurch das Modell End-to-End trainierbar und für anspruchsvolle Anwendungen wie Medizinbildverarbeitung, Robotik oder autonomes Fahren besonders geeignet ist.
Die Entwicklung von Mask R-CNN basiert auf zwei wichtigen Vorläufermodellen, Fast R-CNN und Faster R-CNN, die schrittweise die Grundlage für 
seine heutige Architektur gelegt haben.
Fast R-CNN wurde 2015 von Ross Girshick entwickelt und war eine effiziente Weiterentwicklung des ursprünglichen R-CNN. Es verarbeitet das 
gesamte Bild zunächst durch ein gemeinsames CNN und wendet anschließend RoI Pooling auf vorgeschlagene Regionen an, 
um Feature-Vektoren zu extrahieren. Diese Regionen stammen jedoch aus einem externen, nicht lernbaren Verfahren wie Selective Search, was die 
Geschwindigkeit begrenzt. Das Modell klassifiziert jede Region und optimiert gleichzeitig deren Bounding Box.
Faster R-CNN wurde 2016 von Shaoqing Ren et al. vorgestellt und löst das Geschwindigkeitsproblem von Fast R-CNN durch die Einführung eines lernbasierten RPN. 
Dieses RPN generiert Regionsvorschläge direkt aus dem Feature-Map, wodurch das gesamte Modell vollständig  end-to-end trainierbar wird. 
Die restliche Architektur bleibt ähnlich: RoI Pooling erzeugt Merkmalsvektoren, die zur Klassifikation und Lokalisierung verwendet werden. 
Faster R-CNN bildet somit die direkte Grundlage für Mask R-CNN, liefert jedoch keine Segmentierungsmaske.
Durch die kontinuierliche Weiterentwicklung dieser Modellreihe konnte Mask R-CNN schließlich eine leistungsfähige Lösung für die präzise Instanzsegmentierung 
realisieren, die klassische Objekterkennungsverfahren deutlich übertrifft~\cite{he2017maskrcnn}.

\subsection{U-Net}
U-Net ist ein vollständig konvolutionelles neuronales Netzwerk, das 2015 von Olaf Ronneberger et al. speziell für die 
semantische Segmentierung entwickelt wurde. Es wurde ursprünglich für medizinische Bilddaten konzipiert, eignet sich 
jedoch auch hervorragend für technische Anwendungen, bei denen eine präzise Pixel-zu-Pixel-Zuordnung erforderlich ist – 
etwa in der industriellen Objekterkennung, Robotik oder Visual Servoing. Der Name „U-Net“ stammt von der charakteristischen
U-förmigen Architektur, die aus einem symmetrischen Encoder-Decoder-Aufbau besteht. Der Encoder-Pfad extrahiert hierarchisch Bildmerkmale durch aufeinanderfolgende Convolution- und Max-Pooling-Schritte, wobei die 
räumliche Auflösung reduziert und gleichzeitig die semantische Tiefe der Feature-Maps erhöht wird. Der Decoder-Pfad rekonstruiert die Bildauflösung durch transponierte Faltungen („Upconvolutions“) und kombiniert diese 
mit den hochauflösenden Features aus dem Encoder über Skip-Verbindungen. Diese Verbindungen sorgen dafür, dass feine Details 
aus den frühen Layern nicht verloren gehen – ein entscheidender Vorteil gegenüber früheren Segmentierungsmodellen. 
Das Netzwerk liefert als Ausgabe eine Segmentierungsmaske, in der jedem Pixel eine Klasse zugewiesen wird. U-Net kann auch auf 
kleinen Datensätzen effektiv trainiert werden, da es durch Datenaugmentation und die symmetrische Architektur effizient lernt. 
In der Praxis wird U-Net häufig verwendet, um bestimmte Objektbereiche zu erkennen (z.B. Schrauben, Leitungen), die dann 
gezielt zur Feature-Extraktion oder Steuerung in Vision-Systemen genutzt werden können~\cite{ronneberger2015unet}. 

\subsection{DeepLabV3+}
DeepLabV3+ ist ein leistungsstarkes neuronales Netzwerk zur semantischen Segmentierung, das 2018 von Liang-Chieh Chen et al.
vorgestellt wurde. Es erweitert die vorherige DeepLab-Reihe und kombiniert präzise Lokalisierungsfähigkeiten mit der Fähigkeit, 
kontextuelle Informationen auf mehreren Skalen zu erfassen. DeepLabV3+ ist besonders gut darin, Objekte mit komplexen Formen, 
variablen Größen und feinen Kanten pixelgenau zu segmentieren – was es ideal für Anwendungen in der autonomen Navigation, 
Robotik und Industrievision macht. Das Herzstück von DeepLabV3+ ist der sogenannte Atrous Spatial Pyramid Pooling (ASPP)-Block. 
Dieser verwendet dilated (atrous) Convolutions, um Informationen in verschiedenen räumlichen Auflösungen zu erfassen, ohne die 
Auflösung der Feature-Maps zu stark zu reduzieren. Das bedeutet: Das Netzwerk kann große Objekte vollständig erfassen und 
gleichzeitig feine Details erhalten – ein Balanceakt, den viele klassische CNNs nicht leisten können. Im Gegensatz zu seinem 
Vorgänger DeepLabV3 fügt DeepLabV3+ einen Decoder-Pfad hinzu, der die Segmentierungsergebnisse räumlich verfeinert. Dieser Teil 
des Netzwerks verbessert die Kantenschärfe und Detailgenauigkeit der Segmentierung deutlich – ein entscheidender Vorteil in 
Anwendungsbereichen, wo millimetergenaue Trennung von Objekten gefragt ist. Als Backbone für die Feature-Extraktion kann 
DeepLabV3+ auf leistungsstarke Netzwerke wie Xception oder ResNet zurückgreifen, was die Genauigkeit weiter erhöht. 
Die Architektur ist vollständig konvolutional und eignet sich daher gut für Bilder beliebiger Größe. Sie ist zudem modular 
und flexibel, was sie zu einem beliebten Basismodell in vielen Forschungsarbeiten macht. DeepLabV3+ ist besonders nützlich, 
wenn man Segmentierungsmasken für nachgelagerte Aufgaben wie Merkmalsextraktion oder Objektverfolgung benötigt~\cite{chen2018deeplabv3plus}. 



%% End of File %%